% !TEX program = xelatex
\documentclass[11pt,a4paper]{article}
\usepackage{algo-preamble}

\title{\textbf{L05 \textperiodcentered{} Διαίρει \& Κυρίευε III}\\[2pt]
\large Closest Pair of Points}
\author{Algorithms Class Hub}
\date{\small Αλγόριθμοι και Πολυπλοκότητα (K17) \textperiodcentered{} ΕΚΠΑ}

\begin{document}
\maketitle

\begin{center}\itshape\small
Γεωμετρικό πρόβλημα: βρες το κοντινότερο ζευγάρι σημείων στο επίπεδο. Από
$O(n^2)$ πέφτουμε σε $O(n \log n)$ με διαίρει και κυρίευε και το έξυπνο τέχνασμα
της ζώνης (strip).
\end{center}

\section{Το πρόβλημα}

\textbf{Είσοδος:} $n$ σημεία $p_1, p_2, \dots, p_n$ στο επίπεδο.

\textbf{Ζητούμενο:} βρες ένα \textbf{ζευγάρι} σημείων με τη \textbf{μικρότερη
Ευκλείδεια απόσταση} μεταξύ τους. Δηλαδή ένα $(p_i, p_j)$ που ελαχιστοποιεί
\[
d(p_i, p_j) = \sqrt{(x_i - x_j)^2 + (y_i - y_j)^2}
\]

\subsection{Γιατί μας νοιάζει}

Είναι \textbf{θεμελιώδες γεωμετρικό πρόβλημα}. Εμφανίζεται σε πολλά πραγματικά
συστήματα:
\begin{itemize}
  \item \textbf{Γραφικά υπολογιστών} --- εντοπισμός collisions σε scenes.
  \item \textbf{Όραση υπολογιστών} --- clustering pixels.
  \item \textbf{Συστήματα γεωγραφικών πληροφοριών (GIS)} --- εύρεση πλησιέστερου
  ορόσημου.
  \item \textbf{Σχεδίαση μορίων (computational chemistry)} --- αλληλεπιδράσεις
  ατόμων.
  \item \textbf{Έλεγχος εναέριας κυκλοφορίας} --- εντοπισμός κινδύνου σύγκρουσης
  μεταξύ αεροσκαφών.
\end{itemize}

\section{Αρχικές παρατηρήσεις}

\subsection{Ωμή βία: $\Theta(n^2)$}

Συγκρίνουμε όλα τα ζεύγη --- $\binom{n}{2} = \Theta(n^2)$ συγκρίσεις, κρατάμε το
ελάχιστο. Σωστό αλλά αργό.

\subsection{Μία διάσταση: $O(n \log n)$}

Αν όλα τα σημεία ήταν πάνω σε μια \textbf{ευθεία} (μία συντεταγμένη), το πρόβλημα
θα ήταν εύκολο:
\begin{enumerate}
  \item Ταξινόμησε τα σημεία ως προς $x$ $\to O(n \log n)$.
  \item Σε ένα γραμμικό πέρασμα, βρες το ελάχιστο της διαφοράς $x_{i+1} - x_i$
  $\to O(n)$.
\end{enumerate}
Συνολικά $O(n \log n)$.

\textbf{Διαισθητικά:} σε \textbf{μία} διάσταση, το κοντινότερο ζευγάρι είναι
\textbf{γειτονικό} στην ταξινόμηση --- γιατί αν δεν ήταν, θα υπήρχε κάτι στη μέση
πιο κοντά. Σε δύο διαστάσεις αυτό \textbf{δεν} ισχύει: δύο σημεία μπορεί να έχουν
ίδιο $x$ αλλά να απέχουν πολύ στον $y$.

\begin{keypoint}
\textbf{Στόχος της διάλεξης:} να επιτύχουμε το ίδιο $O(n \log n)$ και σε
\textbf{2 διαστάσεις}. Πρώτη δοκιμή με D\&C δίνει $O(n \log^2 n)$ --- και θα δούμε
ένα ωραίο τρικ που μας πάει στο $O(n \log n)$.
\end{keypoint}

\textbf{Παραδοχή υπόθεσης (για ευκολία):} δύο σημεία \textbf{δεν} μοιράζονται την
ίδια συντεταγμένη $x$. (Αν συμβεί στην εξέταση μπορούμε να τη χαλάσουμε με
ελάχιστη μετατόπιση --- δεν αλλάζει την ασυμπτωτική.)

\section{Πρώτη απόπειρα διαίρεσης --- αποτυγχάνει}

\textbf{Ιδέα:} χωρίζουμε το επίπεδο σε \textbf{4 τεταρτημόρια} και λύνουμε
αναδρομικά σε κάθε ένα. \textbf{Πρόβλημα:} δεν μπορούμε να εγγυηθούμε ότι κάθε
τεταρτημόριο θα έχει περίπου $n/4$ σημεία. Αν όλα τα σημεία τύχουν στο ένα
τεταρτημόριο, η αναδρομή δεν μειώνει τίποτα.

\begin{warningbox}
\textbf{Μάθημα:} σε D\&C χρειάζεται \textbf{εγγύηση} ότι κάθε υποπρόβλημα είναι
αυστηρά μικρότερο από το αρχικό. Στατικά τεταρτημόρια \textbf{δεν} δίνουν αυτή
την εγγύηση. Πρέπει ο διαχωρισμός να είναι \textbf{adaptive} --- να εξαρτάται από
τα δεδομένα.
\end{warningbox}

\section{Σωστή διαίρεση: κάθετη γραμμή στη διάμεσο}

Σχεδιάζουμε \textbf{κάθετη γραμμή} $L$ στην \textbf{κατά $x$ διάμεσο} των
σημείων. Έτσι \textbf{εγγυημένα} $\lfloor n/2 \rfloor$ σημεία είναι αριστερά και
$\lceil n/2 \rceil$ δεξιά.

\textbf{Διαίρει:} βρες την κατά $x$ διάμεσο σε $O(n)$ (θα δούμε πώς αργότερα,
στις δομές δεδομένων) --- γραμμή $L$ $\to$ αριστερό μισό $P_L$, δεξί μισό $P_R$.

\textbf{Κυρίευσε:} αναδρομικά $\delta_1 = \text{ClosestPair}(P_L)$,
$\delta_2 = \text{ClosestPair}(P_R)$, και $\delta = \min(\delta_1, \delta_2)$.

\textbf{Συνδύαζε:} το $\delta$ είναι η μικρότερη απόσταση \textbf{μέσα} σε κάθε
πλευρά. Αλλά μπορεί να υπάρχει ζευγάρι με το ένα σημείο αριστερά και το άλλο
δεξιά που να είναι κοντινότερο. Πρέπει να ελέγξουμε αυτά τα \textbf{μικτά}
ζευγάρια.

\section{Το κρίσιμο τρικ: ζώνη (strip) πλάτους $2\delta$}

Παρατήρηση: αν υπάρχει μικτό ζευγάρι με απόσταση $< \delta$, \textbf{και τα δύο}
σημεία πρέπει να είναι σε απόσταση $< \delta$ από τη γραμμή $L$. Γιατί; Αν ένα
σημείο απείχε $\ge \delta$ από το $L$, τότε ολόκληρη η απόστασή του από
οποιοδήποτε σημείο από την άλλη πλευρά θα ήταν $\ge \delta$ --- δεν μας
ενδιαφέρει.

Άρα \textbf{κόβουμε} όλα τα σημεία που απέχουν περισσότερο από $\delta$ από το
$L$. Μένει μια \textbf{ζώνη πλάτους $2\delta$} γύρω από το $L$.

\textbf{Στη χείριστη περίπτωση} όλα τα σημεία μπορεί να είναι στη ζώνη. Δεν
μπορούμε λοιπόν να βασιστούμε σε «λίγα σημεία στη ζώνη». Χρειαζόμαστε καλύτερη
ιδέα.

\subsection{Ταξινόμηση κατά $y$}

Ταξινομούμε τα σημεία της ζώνης κατά τη συντεταγμένη $y$. Θα πάρει $O(k \log k)$
όπου $k$ το πλήθος των σημείων στη ζώνη (στη χείριστη περίπτωση $k = n$).

\textbf{Ισχυρισμός-κλειδί:} για κάθε σημείο στη ζώνη, \textbf{αρκεί να ελέγξουμε
τις αποστάσεις από τα επόμενα 11 σημεία} στην ταξινόμηση κατά $y$ (κάποιοι
παρουσιάζουν την ίδια ιδέα με αριθμό 7 ή 15 αντί για 11 --- η ασυμπτωτική δεν
αλλάζει).

\section{Γιατί αρκούν 11 γείτονες;}

\begin{keypoint}
\textbf{Λήμμα:} έστω $s_i$ το σημείο με την $i$-οστή κατά $y$ συντεταγμένη στη
ζώνη. Αν $|i - j| \ge 12$, τότε η απόσταση $d(s_i, s_j) \ge \delta$.
\end{keypoint}

\textbf{Απόδειξη με «καρότσι κουτιών».} Χωρίζουμε τη ζώνη πλάτους $2\delta$ σε
\textbf{κουτιά} $\frac{\delta}{2} \times \frac{\delta}{2}$ (δηλαδή 4 στήλες
κουτιών --- 2 αριστερά του $L$, 2 δεξιά).

\textbf{Βήμα 1 --- κάθε κουτί έχει το πολύ 1 σημείο.} Γιατί; Δύο σημεία στο ίδιο
κουτί απέχουν το πολύ τη διαγώνιο του κουτιού, δηλαδή $\sqrt{(\delta/2)^2 +
(\delta/2)^2} = \delta/\sqrt{2} < \delta$. Αλλά αν είναι στην \textbf{ίδια}
πλευρά του $L$ θα είχαμε βρει αυτή την απόσταση αναδρομικά (αντίφαση με τον
ορισμό του $\delta$). Αν είναι σε \textbf{διαφορετικές} πλευρές, πάλι απέχουν
$< \delta$ --- αλλά τότε \textbf{αυτή} είναι μικρότερη από $\delta$, οπότε θα
έπρεπε να την έχουμε βρει. Άρα \textbf{κανένα κουτί} δεν έχει 2 σημεία.

\textbf{Βήμα 2 --- δύο σημεία που διαφέρουν κατά $\ge 3$ σειρές κουτιών έχουν
απόσταση $\ge \delta$.} Γιατί κάθε σειρά είναι ύψους $\delta/2$, οπότε 3 σειρές
διαφορά $\Rightarrow$ κατακόρυφη απόσταση $\ge \delta$.

\textbf{Βήμα 3 --- σε κάθε ομάδα 3 διαδοχικών σειρών χωράνε το πολύ $3 \times 4 =
12$ σημεία.} Άρα αν $|i - j| \ge 12$, τα $s_i, s_j$ είναι σίγουρα σε
διαφορετικές ομάδες 3 σειρών $\to$ απόστασή τους $\ge \delta$.

\begin{intuition}
\textbf{Διαισθητικά:} η ζώνη πλάτους $2\delta$ είναι ένας στενός διάδρομος.
Επειδή τα σημεία δεν μπορούν να είναι πολύ κοντά μεταξύ τους (αλλιώς θα ήταν
μικρότερα από $\delta$, που έχουμε ήδη βρει), σε κάθε «μικρό κομμάτι» του
διαδρόμου χωράνε λίγα σημεία. Άρα η απόσταση που πρέπει να ελέγξω μπροστά μου
είναι \textbf{σταθερή} (11 σημεία), όχι $O(k)$. Αυτό αλλάζει την πολυπλοκότητα
του combine από $O(k^2)$ σε $O(k)$.
\end{intuition}

\section{Ο αλγόριθμος Closest-Pair}

\begin{codebox}
\begin{verbatim}
Closest-Pair(P = p_1, ..., p_n):
  // βάση
  Εάν n <= 3: επίστρεψε ωμή βία O(1)

  // διαίρει
  Υπολόγισε γραμμή L (κατά x διάμεσος) ώστε |P_L| = |P_R| = n/2
  δ1 = Closest-Pair(P_L)
  δ2 = Closest-Pair(P_R)
  δ  = min(δ1, δ2)

  // συνδύαζε
  S = όλα τα σημεία p με απόσταση < δ από τη L
  Ταξινόμησε το S ως προς y
  Για κάθε σημείο s_i στο S:
    Για j = 1 ως 11:
      αν i+j <= |S|:
        δ = min(δ, d(s_i, s_{i+j}))

  επίστρεψε δ
\end{verbatim}
\end{codebox}
\textbf{Επιστρέφει} την ελάχιστη απόσταση. (Παραλλαγή που επιστρέφει και τα δύο
σημεία είναι τετριμμένη --- αρκεί να θυμόμαστε ποιο ζευγάρι την έδωσε.)

\section{Πολυπλοκότητα --- πρώτη εκδοχή: $O(n \log^2 n)$}

Σε κάθε αναδρομική κλήση:
\begin{itemize}
  \item Εύρεση διαμέσου: $O(n)$ (ή ταξινόμηση $O(n \log n)$).
  \item Δύο αναδρομές: $2 T(n/2)$.
  \item Φιλτράρισμα $+$ ταξινόμηση κατά $y$: $O(n \log n)$.
  \item Σάρωση ζώνης με 11 γείτονες: $O(n)$.
\end{itemize}
Άρα $T(n) = 2 T(n/2) + O(n \log n)$. Το Master Theorem \textbf{δεν} εφαρμόζεται
απευθείας στην απλή μορφή (επειδή το $f(n) = n \log n$ δεν είναι $O(n^d)$). Με τη
γενική μορφή ή με δέντρο αναδρομής:
\[
T(n) = O(n \log^2 n)
\]

\section{Βελτίωση σε $O(n \log n)$}

\textbf{Ιδέα:} η ταξινόμηση κατά $y$ είναι ο μόνος λόγος που έχουμε το επιπλέον
$\log n$. Μπορούμε να την κάνουμε \textbf{μία φορά στην αρχή} αντί για κάθε
αναδρομή.

\textbf{Τι αλλάζει:} στο pre-processing ταξινομούμε τα σημεία σε δύο πίνακες ---
$P_x$ (όλα τα σημεία ταξινομημένα κατά $x$) και $P_y$ (όλα τα σημεία
ταξινομημένα κατά $y$). Κάθε αναδρομική κλήση δέχεται \textbf{και τους δύο}
πίνακες περιορισμένους στο υποπρόβλημά της. Δεν χρειάζεται να ξανα-ταξινομήσει.

Με αυτό το trick, η αναδρομή γίνεται $T(n) = 2 T(n/2) + O(n)$. Master Theorem:
$a = 2,\ b = 2,\ d = 1$, case 2 $\to T(n) = O(n \log n)$. Με το $O(n \log n)$
της αρχικής ταξινόμησης, \textbf{συνολικά $O(n \log n)$}.

\begin{keypoint}
\textbf{Τεχνική γενικής χρησιμότητας:} αν έχεις D\&C με $f(n) = n \log n$, ψάξε
αν μπορείς να μετακινήσεις το βάρος της ταξινόμησης \textbf{έξω} από την
αναδρομή. Είναι κλασικό μοτίβο που εμφανίζεται και σε άλλα γεωμετρικά
προβλήματα.
\end{keypoint}

\section{Άσκηση από το PDF}

Οι διαφάνειες τελειώνουν με μια \textbf{κλασική D\&C άσκηση} που αξίζει να
δουλέψεις μόνος σου:

\begin{quoteblock}
Δίνεται πίνακας $A$ μεγέθους $n$ με την ιδιότητα: υπάρχει θέση $i$ τέτοια ώστε
$A[1] < A[2] < \dots < A[i]$ (γνησίως αύξουσα) και $A[i] > A[i+1] > \dots > A[n]$
(γνησίως φθίνουσα). Να δοθεί αλγόριθμος σε χρόνο $O(\log n)$ που εντοπίζει το $i$
(την «κορυφή» του βουνού).
\end{quoteblock}

\textbf{Υπόδειξη (μην το δεις αν θες να το λύσεις):} δυαδική αναζήτηση. Σε κάθε
βήμα κοιτάς το μεσαίο στοιχείο και τους γείτονές του: αν είναι μεγαλύτερο και από
τους δύο, αυτό \textbf{είναι} η κορυφή. Αν ο αριστερός γείτονας είναι μεγαλύτερος,
η κορυφή είναι αριστερά. Αλλιώς δεξιά. Αναδρομή σε μισό $\to T(n) = T(n/2) +
O(1) = O(\log n)$.

\begin{recapbox}
\textbf{Τι κρατάμε από αυτή τη διάλεξη}
\begin{enumerate}
  \item \textbf{Closest Pair} είναι κλασικό γεωμετρικό πρόβλημα. Brute force:
  $O(n^2)$. D\&C: $O(n \log n)$.
  \item \textbf{Σωστή διαίρεση:} κάθετη γραμμή στη διάμεσο, \textbf{όχι} στατικά
  τεταρτημόρια. Εγγυάται $n/2$ $+$ $n/2$ σπάσιμο.
  \item \textbf{Strip τέχνασμα:} μετά την αναδρομή, ξέρουμε ότι το ζευγάρι που
  μας λείπει (αν υπάρχει) είναι μέσα σε ζώνη πλάτους $2\delta$ γύρω από τη
  γραμμή.
  \item \textbf{Γιατί 11 γείτονες;} Επειδή τα σημεία στη ζώνη απαγορεύεται να
  είναι σε απόσταση $< \delta$ μεταξύ τους, σε \textbf{σταθερό} χώρο της ζώνης
  χωράνε \textbf{σταθερά} σημεία.
  \item \textbf{$O(n \log n)$ αντί για $O(n \log^2 n)$:} προταξινόμηση μία φορά
  στην αρχή $\Rightarrow$ κάθε αναδρομή κάνει combine σε γραμμικό χρόνο.
  \item \textbf{Master Theorem:} αν $T(n) = 2T(n/2) + O(n) \Rightarrow
  O(n \log n)$· αν $T(n) = 2T(n/2) + O(n \log n) \Rightarrow O(n \log^2 n)$.
\end{enumerate}
\end{recapbox}

\lecturefooter
\end{document}
